\section{The actual partial-strip quotient at every polynomial weight}
\label{sec:weighted-partial-strip}

The arithmetic block synthesis of MT-M18 gives a direct extension
of the partial-strip result to the draft's polynomial weights.
Its block restrictions need not extend boundedly to the Hilbert
space. Instead, convolution by the smooth block idempotents does
extend, and supplies the finite-rank projections that the argument
requires. This distinction also retains the boundary jets at their
exact continuity thresholds.

Fix one primitive unitary Hecke function and the entire
pole-removed completion $\Xi$ of MT-M18. Let $\psi\in\mathcal E$
be its inverse Mellin transform and
$\mathcal I=\overline{\psi*\mathcal E}^{\,\mathcal E}$.
Fix $a\in\mathbb R$, $u<v$, and $s\ge0$, and put
\[
 \begin{aligned}
 w_{u,v,s}(x)&=\max(e^{2ux},e^{2vx})(1+x^2)^{s/2},\\
 H&=L^2(\mathbb R,w_{u,v,s}(x)\,dx),\qquad
 I_a=T_a\mathcal I,\qquad T_af(x)=e^{ax}f(x),\\
 J&=\overline{I_a}^{\,H},\qquad Q=H/J,\qquad
 q:H\longrightarrow Q,\qquad U_tf(x)=f(x-t).
 \end{aligned}
 \tag{MA-WP.1}
\]
All constants concern this fixed function and these fixed parameters.
There is no hypothesis that the strip contains every zero, and no
uniform separation of interior zeros from its edges.

For $\lambda=\rho-a$, where $\Xi(\rho)=0$ has multiplicity
$m_\rho$, define
\[
 b_s=\#\{j\in\mathbb Z_{\ge0}:2j<s-1\}
     =\max\left(0,\left\lceil\frac{s-1}{2}\right\rceil\right),
 \tag{MA-WP.2}
\]
and the retained multiplicity
\[
 r_\rho=
 \begin{cases}
 m_\rho,&u<\operatorname{Re}(\rho-a)<v,\\
 \min(m_\rho,b_s),&\operatorname{Re}(\rho-a)\in\{u,v\},\\
 0,&\operatorname{Re}(\rho-a)\notin[u,v].
 \end{cases}
 \tag{MA-WP.3}
\]
Thus $s=3$ retains a boundary value, but not a boundary first
derivative. At $s=4$ it retains derivatives of orders zero and one,
up to the actual zero multiplicity.

\paragraph{MA-WP-001 (the literal weighted closure and its boundary blocks).}
The actual closed ideal in (MA-WP.1) is
\[
 J=\left\{f\in H:
    (\mathcal Mf)^{(j)}(\rho-a)=0
       \text{ for every }\rho\text{ and }0\le j<r_\rho\right\}.
 \tag{MA-WP.4}
\]
Every displayed functional is continuous on $H$; at an edge it
is the continuous Mellin boundary derivative. The actual quotient
has finite-dimensional, mutually annihilating bounded block
projections with densely spanning ranges. In a block its algebraic
module is exactly
\[
             \prod_{\rho\ \text{in the block},\ r_\rho>0}
                         \mathbb C[z]/(z^{r_\rho}),
 \tag{MA-WP.5}
\]
with translation acting by multiplication by
$e^{(\rho-a+z)t}$ on the corresponding factor.

Equivalently, replace the maximum in (MA-WP.1) by
$e^{2ux}+e^{2vx}$. Mellin transform then identifies the space
with the analytic Hardy--Sobolev space on $u<\operatorname{Re}z<v$
whose two boundary functions have Sobolev order $s/2$.
In that literal boundary norm, the closed principal ideal generated
by $\Xi(a+z)$ is precisely the simultaneous vanishing ideal in
(MA-WP.4), with all the stated boundary truncations.

\begin{proof}
We give the block construction first; it will also prove the
operator estimates needed for the trace.

\emph{Bounded convolution and the ideal action.}
The weight gives the explicit translation bound
\[
 \|U_t\|_{H\to H}
 \le 2^{s/4}(1+t^2)^{s/4}\max(e^{ut},e^{vt}).
 \tag{MA-WP.6}
\]
Indeed $\max(e^{2u(x+t)},e^{2v(x+t)})$ is at most
$\max(e^{2ut},e^{2vt})\max(e^{2ux},e^{2vx})$, and
$1+(x+t)^2\le2(1+x^2)(1+t^2)$.
The translations are strongly continuous, first on compact smooth
functions and then by density and the local uniform bound.
Consequently, for $h\in\mathcal E$,
\[
 \|C_h\|_{H\to H}
 \le 2^{s/4}\int |h(t)|(1+t^2)^{s/4}
                         \max(e^{ut},e^{vt})\,dt
 \le C_N p_{N,0}(h)
 \tag{MA-WP.7}
\]
for a fixed sufficiently large integer $N$.

The space $I_a$ is a closed convolution ideal of $\mathcal E$,
since $T_a$ is a continuous convolution-algebra automorphism.
If $h\in I_a$, then $C_h(\mathcal E)\subset I_a$.
The inclusion $\mathcal E\subset H$ is continuous and dense,
so (MA-WP.7) implies $C_h(H)\subset J$. Hence every
$h\in I_a$ acts as zero on the actual quotient $Q$.
All convolutions by $\mathcal E$ preserve $J$ and descend to $Q$.

\emph{The precise MT-M18 data used.}
Retain the finite block algebras $(J_n,\|\cdot\|_n)$, block
restrictions $R_n$, and entire interpolation lifts $\mathcal T_n$
of MT-M18, (M18.1)--(M18.4) and (M18.10)--(M18.13).
Their defining facts are
\[
 \begin{gathered}
 d_n=\dim J_n=O(\log(2+|n|)),\qquad \|1\|_n\le1,\\
 R_m\mathcal T_n=\delta_{mn}\operatorname{id},\qquad
 p_{N,j}(\mathcal T_n b)
       \le C_{N,j}(2+|n|)^{D_{N,j}}\|b\|_n,\\
 f-\sum_n\mathcal T_n(R_nf)\in\mathcal I
             \quad(f\in\mathcal E),
 \end{gathered}
 \tag{MA-WP.8}
\]
where the series converges absolutely in every seminorm of
$\mathcal E$. Empty blocks have zero lift. The norm is the
holomorphic quotient norm of MT-M18, not the norm of individual
Hermite coordinates. Thus the displayed constants have no
inverse zero-separation or edge-distance factor.

Center the lifts and restrictions by
\[
 S_n=T_a\mathcal T_n,
 \qquad r_ng=R_n(T_{-a}g),
 \qquad e_n=S_n(1),
 \qquad h_{n,g}=S_n(r_ng).
 \tag{MA-WP.9}
\]
The symbol $r_n$ is used only on $\mathcal E$; no bounded
extension $H\to J_n$ is assumed. For every $g\in\mathcal E$,
\[
 g-\sum_n h_{n,g}\in I_a,
 \qquad \sum_n h_{n,g}\text{ converges in }\mathcal E.
 \tag{MA-WP.10}
\]

\emph{Actual finite-rank projections.}
Let $P_n=C_{e_n}^{Q}$. The full smooth jet description of
MT-M18 implies
\[
 e_n*e_n-e_n\in I_a,\qquad
 e_n*e_m\in I_a\ (n\ne m).
\]
The preceding ideal-action argument therefore gives
$P_n^2=P_n$ and $P_nP_m=0$ for $n\ne m$ on $Q$.
For $f\in\mathcal E$, its block jets give the exact identity
\[
                 P_nq(f)=q(S_n(r_nf)).
 \tag{MA-WP.11}
\]
Thus $P_nq(\mathcal E)$ lies in the finite-dimensional space
$V_n=q(S_nJ_n)$, of dimension at most $d_n$. This space is
closed. Density of $q(\mathcal E)$ and boundedness of $P_n$
show that $P_n(Q)\subset V_n$. Conversely $P_n$ is the
identity on $V_n$, again by its smooth jets. Hence
$P_n(Q)=V_n$ and $P_n$ is finite rank. Equations
(MA-WP.7)--(MA-WP.9) also give
\[
                    \|P_n\|\le C(2+|n|)^D.
 \tag{MA-WP.12}
\]
The block synthesis in (MA-WP.8), followed by $q$, shows
that the algebraic span of the $V_n$ is dense in $Q$.

For later use let $A_{n,g}=C_{h_{n,g}}^Q$.
On the smooth quotient $h_{n,g}$ is supported in block $n$;
therefore the same bounded extension argument gives
\[
 A_{n,g}=P_nA_{n,g}=A_{n,g}P_n,
 \qquad \operatorname{rank} A_{n,g}\le d_n.
 \tag{MA-WP.13}
\]
For the single $N$ in (MA-WP.7), there is an exponent $D$
such that
\[
 \|A_{n,g}\|_1
 \le d_n\|A_{n,g}\|
 \le C d_n(2+|n|)^D\|r_ng\|_n.
 \tag{MA-WP.14}
\]
The last norm decreases faster than every power for
$g\in\mathcal E$, by MT-M18's restriction estimate.
Thus (MA-WP.10) gives the actual operator identity
\[
                  C_g^Q=\sum_n A_{n,g},
 \qquad\text{with convergence in trace norm.}
 \tag{MA-WP.15}
\]
In particular $g\mapsto C_g^Q$ is continuous
$\mathcal E\to\mathcal S_1(Q)$.

This identity removes a possible completion ambiguity.
If $P_n v=0$ for every $n$, (MA-WP.13)--(MA-WP.15)
give $C_g^Qv=0$ for every $g\in\mathcal E$.
Choose a compact smooth approximate identity of integral
one. Strong continuity in (MA-WP.6) gives $C_g^Qv\to v$,
so $v=0$. Therefore
\[
                         \bigcap_n\ker P_n=\{0\}.
 \tag{MA-WP.16}
\]
Dense block ranges alone would not justify this last assertion;
the approximate-identity argument is required.

\emph{Which finite jet combinations are continuous.}
On compact smooth functions write
\[
        \ell_{\lambda,j}(f)
             =\int x^j e^{\lambda x}f(x)\,dx.
\]
Weighted $L^2$ duality gives
\[
 \ell_{\lambda,j}\in H^*
 \quad\Longleftrightarrow\quad
 \int\frac{|x|^{2j}e^{2\operatorname{Re}\lambda x}}
                    {w_{u,v,s}(x)}\,dx<\infty.
 \tag{MA-WP.17}
\]
This holds for every $j$ in the open strip, for
$2j<s-1$ at an edge, and for no $j$ outside the closed
strip.

There is also no cancellation that allows an excluded
higher jet to survive in a finite linear combination.
To see this explicitly, take a finite sum of such jets
that is continuous on $H$. On its finite-dimensional
formal jet space, precomposition with $U_t$ has, at
$\lambda$, the matrix
\[
 \ell_{\lambda,j}(U_tf)
       =e^{\lambda t}\sum_{i=0}^{j}
                       \binom{j}{i}t^{j-i}\ell_{\lambda,i}(f).
 \tag{MA-WP.18}
\]
Choose one real $t$ such that the finitely many numbers
$e^{\lambda t}$ for distinct $\lambda$ are distinct.
A finite Hermite interpolation polynomial in this bounded
dual translation isolates each generalized eigenspace.
Each resulting one-$\lambda$ component is therefore
continuous on $H$. Its integral kernel is
$p(x)e^{\lambda x}$ for a polynomial $p$.
If $p\ne0$ has degree $j$, then $|p(x)|$ is bounded above
and below by positive constants times $|x|^j$ on each
sufficiently distant real tail. Its dual integrability
is exactly (MA-WP.17) for that highest degree.
Hence every nonzero coefficient in a continuous finite
jet combination is among the permitted jets in
(MA-WP.3). No uniform bound for these finite interpolation
polynomials is needed.

\emph{The actual finite blocks and the full closed kernel.}
Let $L\in V_n^*$. The functional $L P_nq$ is continuous on
$H$. Restricted to $\mathcal E$, (MA-WP.11) expresses it
as a finite combination of the full jets in block $n$.
By the preceding argument it is a combination of precisely
the permitted jets. Conversely every permitted jet in that
block is continuous on $H$, vanishes on $J$, and satisfies
$\ell_{\lambda,j}P_n=\ell_{\lambda,j}$ on the quotient.
The finite Hermite interpolation in $J_n$ shows that these
functionals are linearly independent. Thus $V_n^*$ is
exactly their span. The kernel of $J_n\to V_n$ is
therefore the ideal that discards the other coordinates
and the higher boundary powers. This proves (MA-WP.5)
and its translation action.

Every element of $J$ annihilates the permitted jets.
Conversely, if $f\in H$ annihilates them, the description
of $V_n^*$ gives $P_nq(f)=0$ for every $n$.
Equation (MA-WP.16) gives $q(f)=0$, proving (MA-WP.4).

\emph{The literal Hardy--Sobolev interpretation.}
Let $H_\Sigma$ use the weight
$(e^{2ux}+e^{2vx})(1+x^2)^{s/2}$.
The two norms differ by factors between one and $\sqrt2$,
so they have the same closed ideal $J$ and boundedly
similar quotient convolutions. Put $r=s/2$.
The Mellin boundary functions
$F_\sigma(t)=\mathcal Mf(\sigma+it)$, $\sigma=u,v$,
satisfy, by Plancherel,
\[
 \|f\|_{H_\Sigma}^2
  =\frac1{2\pi}\sum_{\sigma\in\{u,v\}}
       \|(1-\partial_t^2)^{r/2}F_\sigma\|_{L^2(dt)}^2.
 \tag{MA-WP.19}
\]
Conversely the $s=0$ strip transform theorem
(MA-HY.3)--(MA-HY.4) gives a
unique underlying $f$ for an analytic Hardy function;
Sobolev regularity of its two boundary values, by the
same Fourier identity, is exactly $f\in H_\Sigma$.
This proves the asserted Hardy--Sobolev identification.
The continuous edge jets are actual derivatives of the
boundary values, also given by the absolutely convergent
integrals in (MA-WP.17).

Finally $\overline{C_{T_a\psi}H_\Sigma}=J$.
Indeed $\mathcal E$ is dense in $H_\Sigma$, convolution
is bounded, and $I_a$ is its $\mathcal E$-closure on the
dense smooth subspace. Under Mellin transform this is
the closure of multiplication by $\Xi(a+z)$.
Thus (MA-WP.4) is the stated closed principal ideal in
the actual Sobolev norm, not an independently imposed
spectral kernel.
\end{proof}

\paragraph{MA-WP-002 (partial trace with exactly the allowed edge multiplicities).}
For every $g\in\mathcal E$, including every compact smooth
test, the actual quotient convolution is trace class and
\[
        \operatorname{Tr}_Q C_g^Q
          =\sum_{\Xi(\rho)=0}r_\rho\,
                                  \mathcal Mg(\rho-a).
 \tag{MA-WP.20}
\]
The sum is absolutely convergent and the map into trace
class is continuous in the topology of $\mathcal E$.
For the actual quotient translation generator $A_Q$,
a sufficiently large real $\kappa$ and an integer $k$
satisfy
\[
                      (\kappa-A_Q)^{-k}\in\mathcal S_1(Q).
 \tag{MA-WP.21}
\]
One may take any $k>D+2$ in terms of the fixed exponent
in (MA-WP.14), and
$\kappa>\max(v,2-a)+1$.

\begin{proof}
Trace class and continuity were proved in (MA-WP.14)--(MA-WP.15)
before any trace was assigned. On $V_n$, the operator
$A_{n,g}$ is multiplication by the class of $\mathcal Mg$.
In the factor $\mathbb C[z]/(z^{r_\rho})$, its diagonal
entries are all $\mathcal Mg(\rho-a)$, so its trace is
$r_\rho\mathcal Mg(\rho-a)$. It vanishes on $\ker P_n$.
Trace-norm convergence in (MA-WP.15) therefore proves
(MA-WP.20). The unconditional local zero count and rapid
vertical decay also prove the displayed scalar series
absolutely convergent directly.

For the resolvent put $F_k(z)=(\kappa-z)^{-k}$.
It is analytic on the fixed strip containing all the
centered MT-M18 block rectangles. Write $[F_k]_n$ for
the class of $w\mapsto F_k(w-a)$ in the original
block algebra $J_n$. These classes satisfy
\[
             \|[F_k]_n\|_n\le C_k(2+|n|)^{-k}.
\]
Apply the centered interpolation lifts to these classes.
The resulting quotient convolutions have rank at most
$d_n$ and, by (MA-WP.7)--(MA-WP.8), trace norms bounded
by $C_kd_n(2+|n|)^{D-k}$. Their sum is trace class for
$k>D+2$.

On the other hand (MA-WP.6) gives the actual resolvent
power as the operator-norm convergent integral
\[
 (\kappa-A_Q)^{-k}
   =\frac1{(k-1)!}\int_0^\infty
              e^{-\kappa t}t^{k-1}U_t^Q\,dt.
 \tag{MA-WP.22}
\]
On each $V_n$ the finite-dimensional translation formula
in (MA-WP.5) identifies this integral with multiplication
by $F_k$. It therefore agrees there with the trace-class
sum just constructed. The block ranges span densely and
both operators are bounded, so they agree on $Q$.
This proves (MA-WP.21).
\end{proof}

\paragraph{MA-WP-003 (the actual ambient lift and cutoff limit for $s\ge3$).}
In particular take $a=1/2$, $u=-\alpha$, $v=\alpha$,
with arbitrary $\alpha>0$ and $s\ge0$, and let
$P_Z$ be the orthogonal projection onto $J^\perp$ in
the exact space $H_{\alpha,s}$ of the draft.
For $g\in C_c^\infty(\mathbb R)$,
\[
 C_gP_Z\in\mathcal S_1(H_{\alpha,s}),\qquad
 \operatorname{Tr}(C_gP_Z)
       =\sum_\rho r_\rho\,\mathcal Mg(\rho-1/2).
 \tag{MA-WP.23}
\]
If $\chi_R$ are compact smooth cutoffs, uniformly bounded
in supremum norm and tending to one almost everywhere,
then
\[
           \|(M_{\chi_R}-1)C_gP_Z\|_1\longrightarrow0.
 \tag{MA-WP.24}
\]
For $g(0)=0$, the actual geometric-order operators satisfy
\[
 \lim_R\operatorname{Tr}(M_{\chi_R}C_gP_G)
       =-\sum_\rho r_\rho\,\mathcal Mg(\rho-1/2),
 \qquad P_G=1-P_Z.
 \tag{MA-WP.25}
\]
These assertions include the requested $s\ge3$ weights.

\begin{proof}
The operator proof of MA-LIFT-002 needs a trace-class
power of the actual quotient resolvent and the explicit
weight. MA-WP-002 supplies the first premise here without
an interior-margin assumption. The weight is unchanged,
so its adjoint generator is still
$L=w^{-1}\partial_x(w\,\cdot)$, and the integration by
parts on the two half-lines, including the finite cusp
jets in (MA-LIFT.12)--(MA-LIFT.16), gives the same bounded
operator plus finite-rank factorization through
$(\kappa-L)^{-k}P_Z$. That operator is trace class by
(MA-WP.21) and Hilbert adjunction. Thus $C_gP_Z$ is trace
class in the specified order. Trace cyclicity with $P_Z$
and (MA-WP.20) give (MA-WP.23).

Uniform boundedness and strong convergence of the
cutoff multipliers imply (MA-WP.24) by finite-rank
approximation in trace norm. The compactly localized
convolution has its usual diagonal trace
$\operatorname{Tr}(M_{\chi_R}C_g)=g(0)\int\chi_R$;
MA-LIFT-003 proves this for the actual weighted space.
Subtracting $M_{\chi_R}C_gP_Z$ gives (MA-WP.25).
This is a trace limit, not a trace-norm convergence
claim for $M_{\chi_R}C_gP_G$ itself.
\end{proof}

\paragraph{Placement in the requested Grothendieck framework.}
The object in (MA-WP.5) is an actual finite-length local
algebra $\mathbb C[z]/(z^{r_\rho})$ in the completed
spectral quotient. The polynomial weight changes its
visible jet length at an edge; it does not change the
arithmetic multiplicity $m_\rho$ of the original entire
function. In the Grothendieck group of the abelian category
of finite-length $\mathbb C[z]$-modules its class has
length $r_\rho$ at its support point, equal to the
class of $r_\rho$ copies of its residue field. Likewise
multiplication by a test transform $G$ has trace
$r_\rho G(\rho-a)$ and determinant $G(\rho-a)^{r_\rho}$.
The determinant here is formed from exterior powers over
$\mathbb C$ of the underlying finite-dimensional vector
space with its endomorphism.
These scalar invariants retain the multiplicity but
do not distinguish the nilpotent Jordan structure
from its semisimplification. The actual truncated
algebra is retained in the operator construction.
This is the nuclear-trace and determinant placement
of BR-TR-01/02 and F-SM-11/12, with inherited cote~1
scans~079--080 and 111--112. Those scans are framework
locators inherited from the predecessor work, not
newly reread source pages. The argument supplies
neither an arithmetic positivity theorem nor an
$\ell$-adic weight realization of this analytic
quotient.

\paragraph{Exact scope and source boundary.}
The new ingredients are the use of bounded convolution
by the MT-M18 block idempotents, the approximate-identity
argument excluding simultaneous-kernel vectors, and the
finite-block boundary classification. The polynomial
lift bound is exactly (M18.13); its proof uses the
Gaussian interpolation kernel (M18.10) and the matching
Stirling and selected-height bounds (M18.6)--(M18.9).
That complete argument was reread for this extension.
No bounded Hilbert block restriction, individual zero
separation estimate, or independent Hardy--Sobolev
cyclicity theorem is imported.

The conclusion is an unconditional computation of the
actual \emph{partial} quotient and its geometric-order
cutoff trace for the fixed Hecke function. It includes
only the multiplicities in (MA-WP.3). It does not turn
that partial trace into the full Meyer trace if omitted
zero jets contribute. MA-HM-003's obstruction to a full
return intertwiner remains, and the full all-zero
transport retains its stated zero-location premise.
